\documentclass{article}
\usepackage[utf8]{inputenc}
\usepackage[T1]{fontenc}
\usepackage[margin=1in]{geometry}
\usepackage{hyperref}
\begin{document}
\begin{center}
{\Large Encyclopaedia Britannica scan \texttt{page-1327-7d41b1db21d5.jpg}}\\[0.5em]
(no running head)
\end{center}
\section*{Continued from an earlier page}
H U E [ 707 ] H U G

Haifa

\textit{Hier.}
der legant all fill it rises gradually to the surface,

and gradually reverted its heights, thinking ama-

ingly, and flinging up great flows. After attin-

ing its greatest height, it gradually falls till it to

daily disappoint. Building \_life drama\_ and boiling

fipings are frequent in molt parts of the island. In many years they replied to the otherup if one he

naive? The molt capital is that which is called Gray

or Ccyfer, in a plain rifling into final hills, and in the

multi of an amphitheater, bounded by the molt marg-

infeces and various shaped the mountains’s among

which the three-headed Hechs fears pre-eminent. See

TerLANN, N 4——Their there's not confined

to the land, they rife in the very few, and form

feuling fourstairs amidth the waves. Their dithers

from the land is unknown, but she ever vulnerable,

twelve molds off the point of Kickeness, emitting fire

and froke, proves that the observations fires and wa-

ters extend to that space; for those avoid effects safe

from the united fury of their two elements.

HUESCA, on ancient and considerable towns of

Spain, in the kingdom of Granada, fieldsed on a plain in W. Long.

1..45, N. Lat. 32.

HUESCAR, or Guuscar, a town of Spain, in the

kingdom of Granada, feated on a plain in W. Long.

1..45, N. Lat. 32.

HUESNE, or Hussx, a familiardan in the Dul-

tie Box, in the Sound, where was the famous obr-

vatory of Tycho Brake. E. Long 13. 5. Nat.

15.54.

HUET (Peter Daniel), a very learned French wri-

ter, born at Caen in Normandy, on the 8th of Re-

burry 1659. He discovered, from his infantry a

great inclination in the duty of palies literature and

the fencences, and at first applied himself to the law;

but De Carteri principles and Boehart served goes

graphy made him change his duties for thalo of phi-

hophyly, mathematics, the languages, and antiqui-

ries. His timismo is Boeharachism him delicious

of knowing him. He contened a very fact friend-

hip with him, and consomata that fared men to

Sweder. Here Chihlen would have engaged him

her teries; but the, feasible of her incumbent temper,

returned to France. All be brought with him was a

copy of a MSI, of Oregon, which he sacrificed at

Stockholm. He refuelle federal offers from Chitlina

after the billions and went to Rome, and from Ga-

tivers as far freeders. In 1670, the defining age

pointed by the king perception to the daughn, his

major, chief of the Hustic for his collegary with our

title of \_this practice\_ on the prices. It was he dust

formed the plan of the commemoration \_in often Dul.\_

\_pliant\_, and directed the execution. His intentions of

piety determined him to enter into holy orders, which

he did at the war we 44.6. Soon after this, he was per-

fetted by the long, no has abby of January and in

168s was nominated to the hillopric of Sollions, which

he exchange for the free of Armouth. After go-

verning their director ten years, his religious, and any

made about of Pontesny near Caen. His love to his

massive place determined him to the there. But that

finis coming upon him, he retired to Paris, and lodged

among the Jefuits in the \_Maifon Profiffe\_, whom he

had made heys to his library. A fewer dilempper

weakened his body extremely, but was the vivacity of

his genius:' he wrote his own life in a very elegant

flyts) and died in 1721, aged 91. He was a run of

very accessible construction, 1nd of great probity, as

well as immense civilization--The following are the

titles of his principal works. 1. \_De time laungre-\_

\_values, of giving some interceptions\_. 2. An eli-

tion of Origen's comentaries on the holy Scriptures,

in Greek and Latin. 3. A travtile on the origin

of the Romanza. \_Dornedible examples, folio.\_

5. \_Sauffinese Alnetane de conocrinia romania et filais.\_

6. Of the duration of the terrestrial passifies, in

French, 7. Albility of the comments on tangica-

tion of the ancient, which has been raindicated into

English. 8. \_Demonstrate du rival and semi porti-\_

\_constitutes.\_ 9. \_Hetainas.\_ 10. Listin and Greek ver-

fen, 6.

HUGHES (John), an ingenious and polite wri-

ter, was born in 1677. In the earliest parts of his

youish, he continued the difference, poetry, drawing,

and minds, in each of which he by team made a con-

fidercable progress() but followed thoof and his other

fluides only as regarded amendments, except frequent

containment on account of his illust of health. The

land Chanelhoe Cowper made hint forecerty for the

commitons of the peace, not his knowledge, and

illinguished him with singular marks of his clients.

He contained in the farm employment under the end

of Meschelek, and had it to the day of his death,

which happened in 1719, the very light in which his

tragedy, entitled \_The large of Damar, was first ac-\_

ing. He was then ... He realized Pontoredley

dialogues of the dead, Vertact's evolutions of Portu-

gal, and the latter of Abrilad and Dolits. He gave

very accurate edition of species's works, with his

life, a gloffary, and remarks; and wrote several papers

in the Father, Speichere, and Grandians Mr Duno-

combo, who married his father, collected his poems

and days in a vols 12mo, in 735.

HUGHILY, a town of Afai, in the kingdom of

Bangels, feated on the molt velocity branch of the in-

ver changes. It is of large extent, reaching about

two miles along the river, flat, and dress a great trade

all all the commodities of this country, affirming cik

cargos for 50 to Ohio annually, besides what is

though to cartages so the neighboring towns,

Solitpere is brought hitner from Patnia in refels above

50 yards long and five broad. The inhabitants are

chely linking but there are also Deutsche Forge,

ilib, and other Europens. E. Long. 87, 55. N.

Las 24.0.

HUGO CAvery, chief of the third race of the kings

of France, being count of Parra and Orleans; he was

refiled to the chrome for his military vollor and public

resturs in 287. See Fears, n 38.

HUGOONIA, in botany 12, genus of the decandria

order, belonging to the monopolph, clic of plants;

and in the natural method ranking with those of which

the order is doubtful. This corridor a postapatolous;

the fruit in a plant with a printed barrel.

HUGUENOTS, an application given by way of

concept to the referred or Protestant Calvinifts of

France.

The name had its first rife in 1560; but authors are

\section*{Source record}
\begin{itemize}
\item Page ID: \texttt{britannica:309f765b1ecc83e0a7f7762e44b4a927a7555443:train:1327}
\item Scan file: \texttt{page-1327-7d41b1db21d5.jpg}
\item Image SHA-256: \texttt{7d41b1db21d5e291f0cbf58aa92f8820d11602e700d5e337935c65b9d75b6211}
\item Source volume ID: \texttt{149981670}
\item Extraction: nvidia/nemotron-3.5-parse via the shared organizer endpoint.
\item Dataset: National Library of Scotland, Encyclopaedia Britannica Illustrated (CC0 1.0).
\end{itemize}
\end{document}
